% \subsection{Post-phenomenological Experience}
\vspace{1.0em}
\enquote{\textit{Technologies can radically transform the situation, including one's sense of one's body} \parencite{Ihde_2002}.}
\vspace{1.0em}

\subsection{Post-Phenomenology: Technology as Mediation of Experience}

Post-phenomenology offers a foundational philosophical account of why technologies cannot be treated as neutral instruments. Instead, we must treat technology as mediators of "human-world relations". Don Ihde’s \emph{Technology and the Lifeworld} makes this case most clearly: technologies are "non-neutral" because every use or interaction both amplifies and reduces the dimensions of its subsequent experience. These technologies actively transform perception and action, not as passive intermediaries \parencite[p. 75]{Ihde_1990}. One example from his text is that of glass, in the sense of magnification technology. The most frequent anecdote he gives is that of Galileo's Telescope. When one looks through a telescope or wears eyeglasses, vision is technologically reconstituted: the technology both extends perception (magnifying detail, correcting focus) and constrains it (narrowing the field, flattening depth). Ihde terms this the \emph{magnification/reduction structure}, an invariant feature of technological mediation that reveals technologies as constitutively transformative, in direct contrast to being the transparent and neutral conduits many design disciplines may have us think they are.\parencite[pp. 75-76]{Ihde_1990}.

Ihde formalizes these insights into four human-technology relations: embodiment, hermeneutic, alterity, and background. In embodiment relations, technologies are incorporated into bodily experience, becoming experientially “transparent.” Yet even when withdrawn into the background of an experience, they still shape that experience by reorganizing perceptual and bodily possibilities \parencite[pp. 72-74]{Ihde_1990}. Hermeneutic relations, by contrast, entail interpretive mediation: we do not see the world directly but read it through technological artifacts, as with thermometers, spectrographs, or musical notation \parencite[pp. 84-93]{Ihde_1990}. In alterity relations, technologies appear as quasi-others. This "technology-as-other" speaks to an interactive presence that we engage with, negotiate, or resist in our relationships with other technologies, and the world at large \parencite[pp. 97-99]{Ihde_1990}. A conversational AI or even a responsive music system becomes something to which we \textit{relate to}, rather than merely use. Lastly, the background relations refer to technologies that fade from view, yet remain infrastructural (in the most normative of cases). For instance, the heating system or the algorithmic recommendation. Perhaps even the electric hum or room tone of a studio \parencite[pp. 108-112]{Ihde_1990}. These operate ambiently, structuring what is possible without calling attention to themselves. In any case, the key point is that mediation is constitutive: technologies are co-actors in experience, and thus in culture.

However, there exists another more extreme (by Ihde's terms), yet more common (in society) reading on these human-technology relations. Building on Ihde, we might also recognize that marginalized knowledges and uses of technologies often inhabit this "background" position. Their apparent lack of agency arises not from absence of will, but from dominant systems that render them infrastructurally invisible and useless, what Ihde terms “experience junk” \parencite[p. 108]{Ihde_1990}.

The implications for digital audio toolsets and their resulting sound artifacts are potentially profound. A digital audio workstation (DAW) or sound synthesis environment do far more than simply emit sound based on our/user parameters or inputs; they frame the actions and interactions of the overall sonic perception through their very design. For instance, the grid structure of most DAWs privileges certain rhythmic assumptions (often common or 4/4 time), amplifying particular creative gestures while simultaneously reducing (flattening) others. Ihde’s framework reveals that what often passes under the guise of “design choice,” in fact, operate as material-cultural inscription: the embedding of values and logics into the very structure of sound-making. To call these tools “neutral” is to ignore their active role in shaping what can be heard, composed, or imagined.

This recognition raises a central critical question: \emph{whose assumptions, practices, and cultural logics are embedded in our technologies?} Ihde reminds us that technologies are always historically and culturally situated---what he calls their "multistability" \parencite[pp. 144-145]{Ihde_1990}. The same instrument may afford different practices in distinct contexts. With Ihde, we return to the example of Galileo’s telescope. The telescope itself did not determine scientific discovery; rather, within a European epistemological framework, interpretations of its use transformed it into a scientific instrument. The same, or another, magnifying technology in a dissimilar context might serve artistic or navigational purposes instead. Likewise, digital audio tools translate sound through the interpretive frameworks of their makers: the grid, the staff, and the scale system are cultural codifications. What Galileo’s telescope did for sight, Western music notation and studio design do for sound, they mediate experience through specific cultural logics that appear universal only because their interpretive frames have become the most dominant ones.

For this research, Don Ihde’s non-neutrality thesis underpins the methodological move toward \emph{technopoiesis}. If technologies always already mediate experience, then the question becomes more about how to deliberately reconfigure such mediations. Post-phenomenology actualizes the conditions under which cultural assumptions enter into technology, while technopoiesis (as defined by authors Juani Guerra and Svend Ostergaard) \parencites{Guerra_Ostergaard_2017} provides a model for how alternative epistemologies can be materially enacted into new tools. Thus, in the case of this thesis, \emph{Sonic Speculocultural Technopoiesis}, Black sonic epistemologies and rituals of Black speculative tradition can become computational design parameters. In short, Ihde alerts us that technological mediation is inevitable and non-neutral; this dissertation extends that insight to argue that mediation can be critically re-authored to embody different cultural worlds. To bring such relations forward is to practice a kind of critical attunement, and how they might be reworked through embodied, collaborative, or fugitive engagements with technology.

Building on post‐phenomenological accounts of mediation, recent scholarship in sound studies underscores that listening itself is a technologically and culturally conditioned act. As LaBelle observes, 

\begin{quotation}
It is my view that sound, as an intensity that moves objects and bodies into the world, extending their reach and relation, is a force that works to link across singularities. In enabling an animate flow to pass between bodies and things, sound is fundamentally a vibrant matter, one that conducts any number of contacts and conversations –-- the rustlings and stirrings by which listening and being heard takes place. The limits of bodies and things are radically extended through sounded actions, making of them expressive flows open to intersections and overlaps, as well as fragmentations and ruptures. Sound \emph{intensifies} relations by animating their potentiality, exposing the matters and bodies of the world to each other. \parencite[p.61]{LaBelle_2018}.  
\end{quotation}

Sound and sound producing technologies script relations of hearing, attention, and authority.  In this sense, sonic mediation extends Ihde’s notion of embodiment relations into the sociosonic domain, revealing how the very possibility of perception is structured through devices that amplify some worlds while silencing others. Ouzounian’s notion of the “sonic undercommons” makes this politics of audition explicit: acts within "collective models of relational listening" can unsettle dominant acoustemologies and open fugitive zones where alternative ways of hearing and knowing take shape \parencite[pp.17-18]{Ouzounian_2021}. Both authors show that sound functions as a site of negotiation between bodies, technologies, and histories rather than as an immaterial by-product of them.

\newpage \subsection{Technopoiesis: Co-Creation of Technology and Knowledge}

If post-phenomenology alerts us that technologies are never neutral mediators of human experience, technopoiesis demonstrates how human cultures, in turn, shape the emergence and stabilization of those technologies. Guerra and Ostergaard (2017) define technopoiesis as the complex dynamic by which knowledge, culture, and material forms co-evolve through processes of emergence and feedback. They describe four interrelated levels of this process: (1) the level of \emph{exigencies}, where human interaction with the environment creates needs and problems; (2) the level of \emph{diagrammatic representation}, where signs and models externalize ideas and allow for experimentation; (3) the level of \emph{material form}, where prototypes and artifacts give provisional embodiment to these representations; and (4) the level of \emph{technological conventionalization}, where technologies become entrenched in cultural practice as stabilized categories \parencites[237-238]{Guerra_Ostergaard_2017}. Another important point on their structure, higher levels emerge from lower ones while simultaneously constraining them, producing a non-linear system of cultural-technical development.

This “dynamicist” view reorients invention away from linear models. While often attributed to isolated discoveries, technologies actually emerge from recursive feedback processes where representations, material experimentation, and cultural stabilization continuously shape one another. For instance, the authors discuss how diagrammatic reasoning often begins as fiction: Leonardo’s designs for flying machines or Maxwell’s electromagnetic wave equations existed first as imaginative constructs before being realized materially \parencites[pp. 249–251]{Guerra_Ostergaard_2017}. In their framework, fiction functions as an essential component of technopoietic processes, opening new trajectories of possible technical futures through speculative diagrams. In this sense, the boundary between cultural imagination and technical reality is porous; speculative representations are themselves a form of technical practice.

Three central concepts, or principles, govern this process: \emph{emergence} (new cultural-technical forms arise from collective interaction without being reducible to individual inventors), \emph{feedback} (stabilized technologies reshape the conditions for further invention), and \emph{synergy} (different knowledge systems co-evolve by constraining and enabling one another) \parencites[251]{Guerra_Ostergaard_2017}. Together these dynamics generate a phenomena that shows, once a technology becomes culturally entrenched, it produces irreversible pathways of complexity, in what the authors describe as\emph{Biopoetics} \parencites[251-252]{Guerra_Ostergaard_2017}. This underscores the importance of the inputs that shape technopoietic processes; once defaults are established, they constrain future developments.

Within this framework, design becomes a practice of shaping how listening, and thus knowing, occurs. The field of \textit{Sonic Interaction Design} (SID) demonstrates how decisions about interface, tone, gesture, and feedback materialize epistemologies into sonic technological formations. Franinović and Serafin emphasize that sonic feedback is an embodied mode of communication through which users orient themselves in action \parencite{Franinovic_Serafin_2013}. In the context of digital audio, technopoiesis situates DAWs, synthesizers, and sound interfaces as outcomes of these emergent feedback loops. Choices about timbre, temporal response, or spatial diffusion thereby inscribe particular assumptions about cognition, emotion, and the body. A post-phenomenological reading of SID reveals that to design a sound, or at the very least the conditions under which one is heard, is to design the very relation between human and world. These conventions, once stabilized, guide how generations of musicians and producers think about and interact with sound. Recognizing this sonic mediation grounds technopoiesis in auditory as well as material practice: when the act of listening becomes a form of making, technological design itself emerges as a site where cultural histories of hearing are encoded, contested, and transformed.

For this dissertation, the technopoietic model presented opens up methodological space for intervention. If representations and cultural logics feed into the cycle, then alternative epistemologies can be intentionally embedded into technological design and their practical artifact-based outcomes. Black sonic practices and temporality can function as inputs that reshape the material and cultural levels of audio technologies. Now within the speculative space, Afrofuturist fictions and sonic imagination act as the technopoietic “diagrams” from which new audio systems may emerge via embodied critical intervention. By engaging with technopoiesis, the research reframes digital audio tools as sites where epistemic imagination and technical reality co-construct one another, opening the possibility for "fugitive" and situated forms of sonic interaction.

